\section{分类与神经网络入门}

\subsection{命名实体识别：一个分类问题}

词向量的一个重要应用是作为分类器的输入特征。以\textbf{命名实体识别}（NER）为例：

\begin{figure}[H]
\centering
\includegraphics[width=0.95\textwidth]{slides-images/slide-035.jpg}
\caption{分类问题的形式化：训练数据 $\{x_i, y_i\}_{i=1}^N$，$x_i$ 为输入（如词向量），$y_i$ 为类别标签\protect\footnotemark}
\end{figure}
\footnotetext{Slides 第36页。}

例如，判断文本中的 ``Paris'' 是人名还是地名：
\begin{itemize}
    \item ``\textbf{Paris Hilton} is famous'' $\rightarrow$ Paris 是\textbf{人名}
    \item ``I visited \textbf{Paris} last spring'' $\rightarrow$ Paris 是\textbf{地名}
\end{itemize}

这需要看上下文才能判断。方法是取中心词及其两侧若干词的词向量，拼接成一个长向量，作为分类器的输入。

\subsection{从线性分类器到神经网络}

传统的统计机器学习分类器（如逻辑回归、SVM、朴素贝叶斯）大多是\textbf{线性分类器}：它们学习权重 $W$，但输入 $x$ 是固定的，决策边界是线性的。

\begin{importantbox}{神经网络分类器的两大优势}
\begin{enumerate}
    \item \textbf{学习输入表示}：不仅学习分类权重 $W$，还同时学习词的分布式表示（词向量本身也可以被更新），这使得在原始符号空间中，分类器是\textbf{非线性的}
    \item \textbf{非线性决策边界}：通过隐藏层和非线性激活函数，可以表示比线性分类器复杂得多的函数
\end{enumerate}
\end{importantbox}

\subsection{窗口分类器的神经网络实现}

Manning 展示了一个简单但完整的神经网络分类器：

\begin{enumerate}
    \item \textbf{输入层}：将窗口内 5 个词（如 ``museums in Paris are amazing''）的词向量拼接为一个 $5d$ 维向量（若 $d = 100$，则为 500 维）
    \item \textbf{隐藏层}：将拼接向量乘以权重矩阵 $W$（如 $8 \times 500$），加偏置 $b$，通过非线性激活函数（如 sigmoid），得到 8 维隐藏表示
    \item \textbf{输出层}：将隐藏表示乘以另一个权重向量，通过 sigmoid 函数输出概率（如``是否为地名''）
\end{enumerate}

\subsection{神经元与 logistic 回归的类比}

\begin{figure}[H]
\centering
\includegraphics[width=0.95\textwidth]{slides-images/slide-040.jpg}
\caption{一个 logistic 回归单元类似于一个简化的神经元：多个输入加权求和后通过非线性激活函数产生输出\protect\footnotemark}
\end{figure}
\footnotetext{Slides 第41页。}

单个人工神经元的数学表达为：

$$h_{w,b}(x) = f(w^T x + b)$$

其中 $f$ 是非线性激活函数（如 sigmoid）：

$$f(z) = \sigma(z) = \frac{1}{1 + e^{-z}}$$

这和 logistic 回归完全相同。\textbf{神经网络的关键在于}：不是只有一个这样的单元，而是有\textbf{一层}这样的单元，它们各自学习不同的权重，提取不同的特征，然后将输出传递给下一层。

\subsection{非线性激活函数的必要性}

\begin{figure}[H]
\centering
\includegraphics[width=0.95\textwidth]{slides-images/slide-045.jpg}
\caption{非线性激活函数的必要性：没有非线性，多层线性变换可以合并为单层（$W_1 W_2 x = Wx$）；有了非线性，网络可以逼近任意复杂的函数\protect\footnotemark}
\end{figure}
\footnotetext{Slides 第46页。}

\begin{warningbox}{没有非线性，深度毫无意义}
如果所有层都是纯线性变换（不加激活函数），那么无论堆叠多少层，整个网络都等价于一个\textbf{单层线性变换}：$W_1 W_2 x = Wx$。只有引入非线性激活函数，每增加一层才真正增加了模型的表达能力，使其能够逼近越来越复杂的函数。
\end{warningbox}

\subsection{交叉熵损失函数}

在 PyTorch 中训练分类器时，最常用的损失函数是\textbf{交叉熵损失}（Cross-Entropy Loss）：

$$H(p, q) = -\sum_c p(c) \log q(c)$$

\begin{itemize}
    \item $p$：真实分布（通常是 one-hot 向量，正确类别为 1，其余为 0）
    \item $q$：模型预测的概率分布
\end{itemize}

当真实标签是 one-hot 时，交叉熵简化为：

$$H(p, q) = -\log q(\text{correct class})$$

即\textbf{负对数似然}（negative log-likelihood）——我们已经很熟悉的损失函数形式。

\begin{knowledgebox}{交叉熵与负对数似然的关系}
交叉熵损失和负对数似然在分类任务中是\textbf{等价的}。在 PyTorch 中使用 \texttt{CrossEntropyLoss} 时，它已经内置了 softmax 和负对数似然的计算，因此模型输出层不需要再手动加 softmax。
\end{knowledgebox}

\subsection{本章小结}

神经网络分类器相比传统线性分类器有两个关键优势：（1）可以同时学习输入表示和分类权重，实现端到端的训练；（2）通过非线性激活函数，可以学习任意复杂的决策边界。一个简单的前馈神经网络由输入层、隐藏层和输出层组成，其中隐藏层通过非线性变换将输入重新表示为更适合分类的形式。

%% ========== 总结与延伸 ==========

